\documentclass[10pt,a4paper]{amsart}
\usepackage[margin=19mm]{geometry}
\usepackage{amsmath,amssymb,mathtools}
\usepackage[scr=rsfs,cal=euler]{mathalfa}
\usepackage{xfrac}
\usepackage[unicode,hidelinks,bookmarks=false]{hyperref}
\newcommand{\ie}{i.e.\ }
\newcommand{\cf}{cf.\ }
\newcommand{\resp}{respectively\ }
\newcommand{\Subsection}{\subsection}
\providecommand{\nobreakdash}{\nobreak\hbox{-}}
\setlength{\emergencystretch}{2em}
\multlinegap0pt
\usepackage{bm}
\usepackage{bbm}
\usepackage{dsfont}
\usepackage{xspace}
\usepackage{cancel}
\usepackage{mathtools}
\usepackage{amsthm}
\makeatletter
\@ifundefined{theorem}{\newtheorem{theorem}{Theorem}}{}
\@ifundefined{lem}{\newtheorem{lem}[theorem]{Lemma}}{}
\@ifundefined{prop}{\newtheorem{prop}[theorem]{Proposition}}{}
\makeatother
\makeatletter
\newcommand{\D}{{\mathcal D}}
\newcommand{\GL}{\operatorname{GL}}
\newcommand{\Hol}{\operatorname{Hol}}
\newcommand{\Real}{\mathbb{R}}
\newcommand{\bt}{{\boldsymbol{\tau}}}
\renewcommand{\d}{\mathrm{d}}
\expandafter\ifx\csname demo\endcsname\relax
\newenvironment{demo}[1]{\par\smallskip\noindent{\bf #1.}}{\par\smallskip}
\fi
\newcommand{\hol}{\mathfrak{hol}}
\expandafter\ifx\csname proclaim\endcsname\relax
\newenvironment{proclaim}[1]{\par\medskip\noindent{\bf #1.}\it}{\par\smallskip}
\fi
\makeatother
\title[Holonomy in pseudo-Hermitian geometry]{Holonomy in pseudo-Hermitian geometry}
\author{Anton S. Galaev, Thomas Leistner, Felipe Leitner}
\date{Source: arXiv:2510.25598v1 (CC BY 4.0)}
\begin{document}
\maketitle

\noindent\emph{Source and licence.} Holonomy in pseudo-Hermitian geometry. Version arXiv:2510.25598v1 (CC BY 4.0), distributed under the Creative Commons Attribution 4.0 International licence, as recorded in the arXiv licence metadata for this exact version and shown on the abstract page https://arxiv.org/abs/2510.25598v1. Full rights notice: licenses/arxiv-2510.25598-CC-BY-4.0.txt.\par\smallskip

\noindent\textit{Editorial scope.} Counted unit: the Theorem whose source label is \texttt{thHolnabl} in the article (source0.tex lines 867--873, source label \texttt{thHolnabl}), delivered together with the article's own complete original proof of it (source0.tex lines 876--975, one \texttt{proof} environment of 100 lines). No mathematics is written, repaired or re-derived: statement and proof are the article's own text line for line. Required context retained uncounted: thholW (lines 606--607); deftildmu (lines 712--714); propHhatistrivial (lines 793--796). Every definition, display and label of those lines is retained unchanged. Omitted from this package and not redistributed: 1--605, 608--711, 715--792, 797--866, 874--875, 976--2478. Nothing inside a retained excerpt is omitted or altered: every retained body is delivered exactly as the article's lines are written, so no omission receipt of any kind accompanies this package. Citations: the retained excerpts carry no citation command at all, so no bibliography entry of the article is transcribed and no reference is redistributed. Reference adaptation: none was needed and none was made; every reference inside the delivered unit resolves inside it, so no unresolved reference or citation is produced. Printed slips kept literal: no printed word, symbol, hypothesis or number of the counted statement or of its proof was altered. Layout and class: the article's own class is \texttt{amsart} with packages including a4wide, amsmath, enumerate, cancel, amsthm, amsfonts, amssymb, graphicx, color, soul, fontenc, verbatim, euscript; the wrapper is the lane's pinned \texttt{amsart} editorial wrapper at 10pt on A4, which restates the theorem-like environments the retained text needs and adds the article's own macro definitions that the retained text needs (\texttt{\textbackslash D}, \texttt{\textbackslash GL}, \texttt{\textbackslash Hol}, \texttt{\textbackslash Real}, \texttt{\textbackslash bt}, \texttt{\textbackslash d}, \texttt{\textbackslash hol}), with extra wrapper packages (bm, bbm, dsfont, xspace, cancel, mathtools). The delivered numbering is the unit's own. Assets: no retained span contains a figure environment, a graphics-inclusion command, a PDF-page inclusion, a table or a TikZ picture; no image file is copied into this package, the package declares no asset dependency and no image file is redistributed with it. Licence: version arXiv:2510.25598v1 of this article is distributed under the Creative Commons Attribution 4.0 International licence, as recorded in the arXiv licence metadata for this exact version and shown on the abstract page https://arxiv.org/abs/2510.25598v1. A full rights notice is delivered at licenses/arxiv-2510.25598-CC-BY-4.0.txt. Characters: every retained excerpt is pure ASCII, so the delivered engine, XeLaTeX with its default Latin Modern text font, reports no missing character.
\begin{theorem}\label{thholW} Let $(M,\theta,g)$ be a contact sub-Riemannian manifold. Then the horizontal holonomy group $\Hol(\nabla)$ coincides with the holonomy group $\Hol(\nabla^W)$ of the Wagner connection. 
\end{theorem}

\begin{equation}\label{deftildmu}
\tilde \mu(t)=\varphi_{f(t)}(\mu(t)),\quad\text{ 
where}\quad f(t)=-\int_a^t\theta(\dot\mu(r))\d r,\quad  t\in[a,b],\end{equation}

\begin{prop}\label{propHhatistrivial} 
Let $(M,g,\theta)$ be a sub-Riemannian contact structure with complete Reeb vector field and let $\nabla^N$ be a metric extended connection.
If $R^N(\xi,\cdot)=0$, then the group $\widehat{\Hol}_x(\nabla^N)$ is trivial.
\end{prop}

\begin{theorem}\label{thHolnabl}
	Let $(M,\theta,g)$ be a contact sub-Riemannian manifold. Suppose that $\bt$ is Codazzi.  Then one of following conditions holds:
	\begin{itemize}
		\item[1.] $\Hol_x(\nabla)=\Hol_x(\nabla^\bt)$;
		\item[2.]  $\Hol_x(\nabla)\subset\Hol_x(\nabla^\bt)$ is a normal subgroup of co-dimension one.
	\end{itemize} 
\end{theorem}

\begin{proof}  Since the subbundle $\hol(\nabla)\subseteq\hol(\nabla^\bt)$ is $\nabla$-parallel,   $\nabla$ preserves the decomposition
$$
\hol(\nabla^\bt)=\hol(\nabla) \oplus \hol(\nabla)^\bot.$$ 
Let $C$ be the skew part of $N^W$ for the Wagner connection $\nabla^W$.  Since $C=\tfrac{1}{4m} R^\bt(\d \theta^{-1})$, we have that $C\in \hol(\nabla^\bt) $. Let $$C=C^\nabla+C^\bot$$ be the corresponding decomposition of the skew-symmetric part $C$ of $N^W$.
For the following recall that from Theorem~\ref{thholW} we have that  \begin{equation}
\label{holWhol}\Hol(\nabla^W)=\Hol(\nabla).\end{equation}


\begin{lem} It holds that $$\nabla^\bt C^\bot=0,$$ and
	$$\hol(\nabla^\bt)=\hol(\nabla) \oplus \left<C^\bot\right>  $$ is the direct sum of ideals.
\end{lem}

\begin{proof}  Since $\bt$ is Codazzi, it holds $R^\bt(\xi,,X)=0$, so that 
$$R^W(\xi,X)=-\nabla_XC=-\nabla_XC^\nabla-\nabla_XC^\bot,\quad\text{for all } X\in\Gamma(\D).$$
Since $R^W(\xi,X)$ takes values in $\hol(\nabla)$, it holds 
$$\nabla C^\bot=0.$$ This implies that $ C^\perp$ commutes with $\Hol(\nabla)$, and with~(\ref{holWhol})  yields  $\nabla^W C^\bot=0$.
%$$L_\xi C^\bot+[\bt,C^\bot]+[C^\nabla,C^\bot]=0.$$
%Since $\nabla C^\bot=0$  it holds that $[\hol(\nabla),C^\bot]=0$. 
We conclude that
$$\nabla^\bt_\xi C^\bot=(\nabla^W_\xi-C^\nabla-C^\bot)\cdot C^\bot=0.$$
Next, since
$$\nabla^\bt_\xi =\nabla^W_\xi-C^\nabla-C^\bot,$$
from~(\ref{holWhol}) we see that $\nabla^\bt_\xi$ preserves $\Gamma(\hol(\nabla))$. This implies that the bundle $\hol(\nabla)$ is $\nabla^\bt$-parallel.
For each point $y\in M$ and vectors $X,Y\in C_x$ it holds 
$$R^\bt_y(X,Y)=R^W_y(X,Y)+(d\theta)_y(X,Y)C_y\in\hol_y(\nabla)\oplus\Real (C^\bot)_y.$$
We conclude that for each curve $\gamma$ from $x$ to $y$ it holds
$$(\tau^\bt_\gamma)^{-1}\circ R^\bt_y(X,Y)\circ\tau^\bt_\gamma\in  \hol_x(\nabla)\oplus\Real (C^\bot)_x.$$ Note also that $R^\bt(\xi,\cdot)=0.$
Thus, by the Ambrose--Singer Theorem, 
$$\hol_x(\nabla^\bt)=\hol_x(\nabla) \oplus \left<C_x^\bot\right>.$$ \end{proof}

Note that the statement of the above Lemma also holds  in the case $C^\bot=0$.

\begin{lem}\label{lemxi1} Let $\lambda:[0,r]\to M$ be an integral curve of the vector field $\xi$. Then it holds
	$$\tau_\lambda^\bt=e^{r\cdot C^\bot_{\lambda(r)}}\circ \tau_\lambda^W\circ B(r),$$ where $B(r)\in\Hol_{\lambda(0)}(\nabla)$.
\end{lem}

\begin{proof}  For each $s\in [0,r]$, let 
$$B(s)=((\tau_\lambda^W)_s)^{-1}\circ e^{-s\cdot C^\bot_{\lambda(s)}}\circ (\tau_\lambda^\bt)_s:\D_{\lambda(0)}\to \D_{\lambda(0)},$$
be a curve in $\GL(\D_{\lambda(0)})$,
where $$(\tau_\lambda^\bt)_s,(\tau_\lambda^W)_s:\D_{\lambda(0)}\to\D_{\lambda(s)}$$ are parallel transports along the curve $\lambda|_{[0,s]}$.
Let $X(s)$ be a $\nabla^\bt$-parallel section of $\D$ along the curve $\lambda$, i.e., it holds
$$\nabla^\bt_{\dot\lambda(s)}X(s)=0,\quad (\tau_\lambda^\bt)_s X(0)=X(s). $$
Then, using that   $\nabla^W C^\perp=0$, we have 
\begin{eqnarray*}
0&=&\nabla_{\dot\lambda(s)}^\bt X(s)
\\
&=&\left(\nabla^W_{\dot\lambda(s)}-C_{\lambda(s)} \right) \left( e^{s\cdot C^\bot_{\lambda(s)}}\circ (\tau_\lambda^W)_s\circ B(s)\right) X(0)
\\
&=& \left(C^\bot_{\lambda(s)}\circ e^{s\cdot C^\bot_{\lambda(s)}}\circ (\tau_\lambda^W)_s\circ B(s)\right) X(0)+\left( e^{s\cdot C^\bot_{\lambda(s)}}\circ (\tau_\lambda^W)_s\circ B'(s)\right) X(0)
\\
&&{}-\left(C_{\lambda(s)}\circ e^{s\cdot C^\bot_{\lambda(s)}}\circ (\tau_\lambda^W)_s\circ B(s)\right) X(0).
\end{eqnarray*}
With  $[C,C^\perp]=0$, this implies 
$$(\tau_\lambda^W)_s\circ B'(s)-C^\nabla_{\lambda(s)}\circ (\tau_\lambda^W)_s\circ B(s)=0,$$
so that
$$ B'(s)\circ B^{-1}(s)=((\tau_\lambda^W)_s)^{-1}\circ C^\nabla_{\lambda(s)}\circ (\tau_\lambda^W)_s.$$ 
Since $C^\nabla\in \hol(\nabla)$, this  shows that $B(s)$ is the development of the curve 
$((\tau_\lambda^W)_s)^{-1}\circ C^\nabla_{\lambda(s)}\circ (\tau_\lambda^W)_s\in\hol_{\gamma(0)}(\nabla)$. \end{proof}

\begin{lem}\label{lemxi2} Let $\mu:[a,b]\to M$ be an arbitrary curve. Then there exist a horizontal curve $\bar\mu:[a,b]\to M$ such that
	$$\bar\mu(a)=\mu(a),\quad \bar\mu(b)=\mu(b),\quad \tau_{\bar\mu}=\tau^W_\mu.$$ 
\end{lem}

\begin{proof}  Let $\check\mu$ be a horizontal curve connecting the points $\mu(a)$ and $\mu(b)$. We obtain the loop $\mu*\check\mu^{-1}$ at the point $\mu(a)$.
Since the holonomy groups of the connections $\nabla$ and $\nabla^W$ coincide, there exists a horizontal loop $\hat\mu$ at the point $\mu(a)$ such that
$$\tau^W_{\mu*\check\mu^{-1}}=\tau_{\hat\mu}.$$ This implies that 
$\tau^W_{\mu}=\tau_{\hat\mu*\check\mu}$.
\end{proof}


Let $\gamma:[a,b]\to M$ be an arbitrary loop at the point $x$.
Since $\gamma([a,b])\subseteq M$ is compact, there exist
numbers $$a_0=a< a_1<\dots< a_{k-1}<a_k=b$$ such that,
for each restriction
$$\gamma_i=\gamma|_{[a_{i-1},a_i]},$$ the horizontal curve $\tilde \gamma_i$ defined by \eqref{deftildmu} is well-defined. 
%\tcom{Need to explain $\tilde{\gamma}$ here!}\acom{Is it enough to refer to \eqref{deftildmu}?}
By the definition of $\tilde\gamma_i$, there exists an integral curve $\lambda_i(s)$, $s\in [0,r_i]$ of the vector field $\xi$ connecting the end-points $\tilde\gamma(a_i)$ of the curve $\tilde\gamma_i$ with the end-point $\gamma(a_i)$ of the curve $\gamma_i$, and it holds 
% \tcom{$r_i=- \int ...$?}\acom{here with +, because $\lambda_i$ goes from end-points of  $\tilde\gamma_i$ to the end-point of $\gamma_i$ }
$$r_i= \int_{\gamma_i}\theta.$$
Since $R^\bt(\xi,\cdot)=0$,  
%\tcom{Why the implication?} \acom{explained} 
the arguments from the proof of Proposition \ref{propHhatistrivial} imply
$$\tau^\bt_{\gamma_i}=\tau^\bt_{\lambda_i}\circ \tau_{\tilde\gamma_i}.$$
From Lemmas \ref{lemxi1} and \ref{lemxi2} it follows that there exists a horizontal curve $\mu_i$ connecting the points $\gamma(a_{i-1})$ and $\gamma(a_i)$ such that 
$$\tau^\bt_{\gamma_i}=e^{r_iC^\bot_{\gamma(a_i)}}\circ\tau_{\mu_i}.$$
Using this, the equality
$$\tau^\bt_\gamma=\tau^\bt_{\gamma_k}\circ\cdots\circ\tau^\bt_{\gamma_1},$$
and applying repeatedly the fact that $C^\bot$ is $\nabla$-parallel, we get 
$$\tau^\bt_\gamma=e^{(r_1+\cdots +r_k)C^\bot_x}\circ\tau_{\mu},$$
where $\mu$ is a horizontal loop at the point $x$.
It is obvious that 

$$r_1+\cdots +r_k=\int_\gamma\theta.$$
Thus, $$\tau^\bt_\gamma=\exp\left(\left(\int_\gamma\theta\right)\cdot {C^\bot_x}\right)\circ\tau_{\mu}.$$
Since
$C^\perp$ commutes with $\Hol(\nabla)$, this implies that the subgroup $\Hol_x(\nabla)\subseteq\Hol_x(\nabla^\bt)$ is normal. It is obvious that the map
$$\lambda:\Real_+\to\Hol_x(\nabla^\bt)/\Hol_x(\nabla),$$
$$\lambda:\exp\left(\int_\gamma\theta\right)\mapsto
\exp\left(\left(\int_\gamma\theta\right)\cdot {C^\bot_x}\right)\cdot \Hol_x(\nabla)$$
is well-defined and surjective. \end{proof}

\end{document}
